% Generated draft: requires mathematical and source review.
\documentclass[11pt]{article}
\usepackage[margin=25mm]{geometry}
\usepackage{fontspec}
\setmainfont{Latin Modern Roman}
\usepackage{amsmath,amssymb,mathtools,unicode-math}
\setmathfont{Latin Modern Math}
\usepackage{enumitem,needspace,xurl,hyperref}
\setlength{\parindent}{0pt}
\setlength{\parskip}{5pt}
\setlength{\emergencystretch}{3em}
\newcommand{\N}{\mathbb{N}}
\newcommand{\Z}{\mathbb{Z}}
\newcommand{\Q}{\mathbb{Q}}
\newcommand{\R}{\mathbb{R}}
\newcommand{\C}{\mathbb{C}}
\newcommand{\F}{\mathbb{F}}
\newcommand{\A}{\mathbb{A}}
\newcommand{\PP}{\mathbb{P}}
\newcommand{\E}{\mathbb{E}}
\newcommand{\eps}{\varepsilon}
\newcommand{\dd}{\,\mathrm{d}}
\begin{document}
\section*{UnsolvedMath ID: 6200086}
\textbf{Source problem code:} catalogue entry
\subsection{Definitions and mathematical statement}
A group $G$ is coherent if
\[
\forall H\leq G,\qquad
H\text{ finitely generated}\ \Longrightarrow\ H\text{ finitely presented}.
\]
This is the explicit algebraic property recalled in the source. It is a condition on every finitely generated subgroup, not merely on $G$ itself.

The actual problem is the broader request to find new restrictions on Kleinian groups. In a real-hyperbolic interpretation, these are discrete groups of isometries of $\mathbb H^n$; the classical three-dimensional usage concerns discrete subgroups of $\operatorname{PSL}_2(\C)$. The source does not specify the ambient dimension, torsion assumptions, finite generation, or geometric finiteness conditions.

No particular new restriction is proposed, and the recalled coherence definition is not itself a conjecture that all higher-dimensional Kleinian groups are coherent. The supplied summary suggests higher-dimensional questions after three-dimensional progress, but does not determine a single class or quantified assertion. A precise theorem target must first identify both the Kleinian-group class and the proposed algebraic or geometric restriction.
\subsection{Short English statement}
Find additional algebraic or geometric constraints on Kleinian groups. Coherence is supplied as context, without a conjecture asserting it for all dimensions.
\subsection{Sources}
\begin{itemize}
\item Local UnsolvedMath registry, permanent ID 6200086, exact source code catalogue entry, ``Boundaries of Groups and Kleinian Groups---Problem 86''; proposed year 2005, no proposer supplied. Registry status open, classification OPEN-TRIAGE.
\item Kapovich, \emph{Problems on Boundaries of Groups and Kleinian Groups} (2005), Problem 86, PDF page 22, \url{https://www.math.ucdavis.edu/~kapovich/EPR/problems.pdf}. Supplied normalized pdf-text extraction accessed July 29, 2026. No PDF was accessed here.
\item The supplied summary calls higher-dimensional restrictions open and describes three-dimensional coherence as settled, using the label ``Scott conjecture'' and an Agol attribution. It gives no precise theorem or reference; the relationship among compact-core/coherence and tameness statements and their attributions requires checking, so no exact historical theorem is inferred from that sentence.
\item Provenance flags current status and statement-reuse terms NEEDS\_REVIEW; the source is publicly author-hosted.
\end{itemize}
\end{document}
